\subsection{分式环中的理想}

命题 11. 设 A 是环，S 是 A 的乘性子集。对每个理想 \(\mathfrak{b}'\)（\(\mathbf{S}^{-1}\mathbf{A}\) 的），令 \(\mathfrak{b} = (i_{\mathbf{A}}^{\mathbf{S}})^{-1}(\mathfrak{b}')\)，则 \(\mathfrak{b}' = \mathbf{S}^{-1}\mathfrak{b}\)。

\begin{enumerate}
\def\labelenumi{(\roman{enumi})}
\item
  设 \(f\) 是典范同态 \(\mathbf{A} \to \mathbf{A} / \mathfrak{b}\)。从 \(\mathbf{S}^{-1}\mathbf{A}\) 到 \((f(\mathbf{S}))^{-1}(\mathbf{A} / \mathfrak{b})\) 的、与 \(f\) 典范相伴的同态（第 1 节，命题 2）是满射，其核为 \(\mathfrak{b}'\)，由此通过取商定义了 \((\mathbf{S}^{-1}\mathbf{A}) / \mathfrak{b}'\) 到 \((f(\mathbf{S}))^{-1}(\mathbf{A} / \mathfrak{b})\) 上的一个典范同构。此外，从 \(\mathbf{A} / \mathfrak{b}\) 到 \((f(\mathbf{S}))^{-1}(\mathbf{A} / \mathfrak{b})\) 的典范同态是单射。
\item
  映射 \(\mathfrak{b}' \mapsto \mathfrak{b} = (i_{\mathbf{M}}^{\mathbf{S}})^{-1}(\mathfrak{b}')\) 限制到 \(\mathbf{S}^{-1}\mathbf{A}\) 的极大（相应地，素）理想的集合上，是该集合到 \(\mathbf{A}\) 的在那些不与 \(\mathbf{S}\) 相交的理想中为极大的理想组成的集合（相应地，\(\mathbf{A}\) 的不与 \(\mathbf{S}\) 相交的素理想组成的集合）上的一个同构（关于包含关系）。
\item
  若 \(\mathfrak{q}'\) 是 \(S^{-1}A\) 的素理想且 \(\mathfrak{q} = (i_A^S)^{-1}(\mathfrak{q}')\)，则存在分式环 \(A_{\mathfrak{q}}\) 到环 \((S^{-1}A)_{\mathfrak{q}}\) 上的一个同构，它把 \(a/b\) 映为 \((a/1)/(b/1)\)，其中 \(a \in A\)、\(b \in A - \mathfrak{q}\)。

\end{enumerate}

\begin{enumerate}
\def\labelenumi{(\roman{enumi})}
\item
  \((f(S))^{-1}(A/b)\) 可借助这两个模之间的典范同构（第 2 节，命题 6）与 \(S^{-1}(A/b)\) 等同。于是正合列 \(0 \rightarrow b \rightarrow A \rightarrow A/b \rightarrow 0\) 诱导出正合列

\[
0 \rightarrow \mathrm{S} ^ {- 1} \mathfrak {b} \rightarrow \mathrm{S} ^ {- 1} \mathrm{A} \rightarrow \mathrm{S} ^ {- 1} (\mathrm{A} / \mathfrak {b}) \rightarrow 0
\]

（第 4 节，定理 1）；它的存在性证明了 (i) 的第一个断言，这里用到 \(\mathfrak{b}' = \mathrm{S}^{-1}\mathfrak{b}\) 这一事实。由于 \(\mathfrak{b}\) 关于 S 是饱和的，条件 \(a \in \mathbf{A}\)、\(s \in \mathbf{S}\)、\(as \in \mathfrak{b}\) 蕴含 \(a \in \mathfrak{b}\)；于是比为 \(s\) 的位似映射（在 \(\mathbf{A}/\mathfrak{b}\) 上）是单射，这就证明了 (i) 的第二个断言。

\item
  我们首先指出，关系 \(b' = S^{-1}A\) 等价于关系 \(b \cap S \neq \varnothing\)，后者表达了 \(b'\) 含 \(S^{-1}A\) 的可逆元这一事实。由第 4 节命题 10 (iii) 可知，\(b' \mapsto b = (i_{\mathrm{A}}^{\mathrm{S}})^{-1}(b')\) 是 \(S^{-1}A\) 的异于 \(S^{-1}A\) 的理想组成的集合到 A 的不与 S 相交且满足命题 10 的条件 (MS) 的理想组成的集合 F 上的一个同构（关于包含关系）。若 \(b'\) 是极大的（相应地，素的），显然 \(b'\) 在 F 中极大（相应地，是素的），反之亦然（由 (i)）。另一方面，若 r 是 A 的与 S 不相交的理想，则它的饱和 \(r_{1}\)（关于 S 的）是 A 的一个包含 r 且与 S 不相交的理想：因为没有元素 \(a \in S\) 能满足 \(sa \in r\)（对某个 \(s \in S\)），否则将得出 \(sa \in r \cap S\)。由此得出，若 r 在 A 的与 S 相交的理想当中为极大的，则它在 F 中极大。类似地，若 r 是不与 S 相交的素理想，则按素理想的定义它满足第 4 节命题 10 的条件 (MS)，从而属于 F。这就完成了 (ii) 的证明。
\item
  设 \(q'\) 是素的并且 q 也是素的。集合 \(T = A - q\) 是 A 的一个包含 S 的乘性子集，从而 \(ST = T\)。我们记 \(T' = i_{\mathbf{A}}^{\mathbf{S}}(\mathbf{T})\)；由第 3 节命题 7 (i) 可知，存在唯一的同构 j，从 \(T^{-1}A = A_{q}\) 到 \(T'^{-1}(S^{-1}A)\)，使得

\[
j (a / b) = (a / 1) / (b / 1),
\]

其中 \(a \in \mathbf{A}\)、\(b \in \mathbf{T}\)。另一方面，\(\mathbf{T}'\) 显然不与 \(\mathfrak{q}'\) 相交；反之，设 \(a / s \in \mathrm{S}^{-1}\mathrm{A}\)；由于 \(1 / s\) 在 \(\mathrm{S}^{-1}\mathrm{A}\) 中可逆，条件 \(a / s \notin \mathfrak{q}'\) 等价于 \(i_{\mathbf{A}}^{\mathbf{S}}(a) = a / 1 \notin \mathfrak{q}'\)，从而等价于 \(a \notin \mathfrak{q}\)；由此得出 \(\mathrm{S}^{-1}\mathrm{A} - \mathfrak{q}' = \mathrm{S}^{-1}\mathrm{T}'\)，于是由第 3 节命题 8 得 \(\mathrm{T}'^{-1}(\mathrm{S}^{-1}\mathrm{A}) = (\mathrm{S}^{-1}\mathrm{A})_{\mathfrak{q}'}\)。

(iii) 中定义的同构称为典范同构。若 A 是整环，则 \(A_{q}\) 与 \((\mathrm{S}^{-1}\mathrm{A})_{q'}\) 到 A 的分式域 K 的子环上的典范同构具有相同的象。
\end{enumerate}

注. 为使 A 的理想 a 满足 \(S^{-1}a = S^{-1}A\)（或者，由第 4 节定理 1，这等价于 \(S^{-1}(A/a) = 0\)），必须且只需 \(a \cap S \neq \varnothing\)，这由定义立即得出。

推论 1. 设 A 是环，S 是 A 的乘性子集。A 的在不与 S 相交的理想中为极大的每个理想 p 都是素理想。

由命题 11，关于 p 的假设意味着 \(\mathfrak{p} = (i_{\mathrm{A}}^{\mathrm{S}})^{-1}(\mathfrak{m}')\)，其中 \(m'\) 是 \(S^{-1}A\) 的极大理想；由于 \(m'\) 是素的，p 也是素的。

推论 2. 设 A 是环，S 是 A 的乘性子集。对 A 的每个不与 S 相交的理想 a，存在一个包含 a 且不与 S 相交的素理想。

\(S^{-1}a \neq S^{-1}A\)（注），因而存在 \(S^{-1}A\) 的一个包含 \(S^{-1}a\) 的极大理想（《代数学》第 I 章，§ 8，第 7 节，定理 2），于是由命题 11 (ii) 即得本推论。

推论 3. 设 A、B 是两个环，\(\rho\) 是 A 到 B 的同态，\(\mathfrak{p}\) 是 A 的素理想。为使存在 B 的素理想 \(\mathfrak{p}'\) 满足 \(\overline{\rho}^{1}(\mathfrak{p}') = \mathfrak{p}\)，必须且只需 \(\overline{\rho}^{1}(\mathrm{B}\rho(\mathfrak{p})) = \mathfrak{p}\)。

若存在理想 \(p'\)（\(B\) 的）使得 \(\overline{\rho}^{1}(p') = p\)，则 \(\rho(p) \subset p'\)，从而 \(B\rho(p) \subset p'\)，且 \(\overline{\rho}^{1}(B\rho(p)) \subset \overline{\rho}^{1}(p') \subset p\)；由于反向包含是明显的，故 \(\overline{\sigma}^{1}(B\rho(p)) = p\)。反之，设 \(\overline{\rho}^{1}(B\rho(p)) = p\)，并考虑乘性子集 \(S = \rho(A - p)\)（\(B\) 的）；假设表明

\[
\mathrm{S} \cap \mathrm{B} _ {\rho} (\mathfrak {p}) = \varnothing ;
\]

由推论 2，存在一个素理想 \(\mathfrak{p}'\)（B 的），它包含 \(\mathrm{B}\varphi(\mathfrak{p})\) 且不与 S 相交；于是 \(\overline{\rho}^{1}(\mathfrak{p}')\) 包含 \(\mathfrak{p}\)，且不能包含 A - \(\mathfrak{p}\) 的任何元素，从而等于 \(\mathfrak{p}\)。

推论 4. 设 A、B 是两个环，\(\rho\) 是 A 到 B 的同态。

\begin{enumerate}
\def\labelenumi{(\roman{enumi})}
\item
  假设存在一个 B-模 E，使得 \(\rho_{*}(\mathbf{E})\) 是忠实平坦的 A-模。则对 A 的每个素理想 \(\mathfrak{p}\)，存在 B 的一个素理想 \(\mathfrak{p}'\)，使得 \(\overline{\rho}^{-1}(\mathfrak{p}') = \mathfrak{p}\)。
\item
  反之，假设 B 是平坦的 A-模。这时，若对 A 的每个素理想 p 存在理想 \(p'\)（B 的）使得 \(\overline{\rho}^{-1}(p') = p\)，则 B 是忠实平坦的 A-模。
\end{enumerate}

\begin{enumerate}
\def\labelenumi{(\roman{enumi})}
\item
  该假设蕴含：对每个理想 \(\mathfrak{a}\)（\(\mathbf{A}\) 的），\(\bar{\rho}^{1}(\mathrm{B}\varphi (\mathfrak{a})) = \mathfrak{a}\)（第 I 章，§ 3，第 5 节，命题 8 (ii)）；于是只需应用推论 3。

\item
  只需证明：对 A 的每个极大理想 m，存在 B 的一个极大理想 \(m'\) 使得 \(\bar{\rho}^{-1}(m') = m\)（第 I 章，§ 2，第 5 节，命题 9 (e)）。现在由假设存在 B 的理想 q 使得 \(\bar{\rho}^{-1}(q) = m\)；由于 \(q \neq B\)，存在 B 的一个包含 q 的极大理想 \(m'\)，从而 \(\bar{\rho}^{-1}(m') \supset m\)；但由于 \(\bar{\rho}^{-1}(m')\) 不能包含 1，故 \(\bar{\rho}^{-1}(m') = m\)。
\end{enumerate}

推论 5. 设 A 是环，S 是 A 的乘性子集，B 是满足 \(i_{\mathbf{A}}^{\mathbf{S}}(\mathbf{A}) \subset \mathbf{B} \subset \mathbf{S}^{-1}\mathbf{A}\) 的环。设 \(\mathfrak{q}\) 是 B 的素理想，使得 A 的素理想 \(\mathfrak{p} = (i_{\mathbf{A}}^{\mathbf{S}})^{-1}(\mathfrak{q})\) 不与 S 相交，并设 \(\mathfrak{p}'\) 是素理想 \(\mathbf{S}^{-1}\mathfrak{p}\)（\(\mathbf{S}^{-1}\mathbf{A}\) 中的）。则 \(\mathfrak{p}' \cap \mathbf{B} = \mathfrak{q}\)。

设 \(S' = i_{\mathbf{A}}^{\mathbf{S}}(\mathbf{S})\)；从 \(S'^{-1}\mathbf{B}\) 到 \(S^{-1}\mathbf{A}\) 的一个典范同构已在第 1 节（命题 2 的推论 4）中定义；我们借助这一同构把这两个环等同起来。由于 \(q \cap S' = \varnothing\)，\(q' = S'^{-1}q\) 是 \(S^{-1}\mathbf{A} = S'^{-1}\mathbf{B}\) 的唯一素理想，使得 \(q' \cap B = (i_{\mathbf{A}}^{\mathbf{S}})^{-1}(q') = q\)（命题 11 (ii)），从而 \((i_{\mathbf{A}}^{\mathbf{S}})^{-1}(q') = p\)；于是 \(q' = p'\)（命题 11 (ii)）。

在推论 5 的记号下，\(\mathbf{A}_{\mathfrak{p}}\) 与 \(\mathbf{B}_{\mathfrak{q}}\) 到 \((\mathrm{S}^{-1}\mathrm{A})_{\mathfrak{p}}\) 上有典范同构（命题 11 (iii)），从而有一个典范同构 \(\mathbf{A}_{\mathfrak{p}} \to \mathbf{B}_{\mathfrak{q}}\)。

